\section{Công Trình Liên Quan}\label{sec:related-work}

Phần này trình bày công trình liên quan chia thành năm nhóm: (i) suy luận tĩnh trên hình thức luận
chung, (ii) ngữ nghĩa vận hành theo vết thực thi, (iii) căn chỉnh định lượng và liên tục, (iv) đối
chiếu với ràng buộc phẳng bên ngoài, và (v) đặc tả và kiểm chứng cấu trúc vai trò tổ chức. Bốn nhóm
đầu giải quyết bài toán kiểm tra tính nhất quán giữa mô hình mục tiêu và mô hình tiến trình đã tồn
tại độc lập, đúng phạm vi của RQ2; nhóm thứ năm đặc tả và kiểm chứng cấu trúc tổ chức tách biệt khỏi
cả mô hình mục tiêu lẫn mô hình tiến trình, đúng phạm vi của RQ1. Mỗi nhóm được trình bày theo cùng
một cấu trúc: cơ chế chung, các đại diện tiêu biểu, và điểm khác biệt cụ thể giữa phương pháp đề
xuất với nhóm đó.

\subsection{Suy Luận Tĩnh Trên Hình Thức Luận Chung}\label{sec:rw-static}

Nhóm công trình đầu tiên dịch cả mô hình mục tiêu và mô hình tiến trình sang một hình thức luận
chung rồi dùng một bộ máy suy luận hoặc truy vấn tự động để phân loại tính nhất quán, không mô phỏng
thực thi. Ghose và Koliadis~\cite{ghose2007validation} dịch goal model GRL và process model BPMN
sang Description Logics, rồi dùng reasoner Pellet để phân loại từng cặp workflow pattern và
intentional relation thành ba trạng thái: realization equivalent, khi mọi đường thực thi đều thoả
goal; strong inconsistency, khi không đường thực thi nào thoả; và potential inconsistency, khi chỉ
một số đường thực thi thoả. Gröner và cộng sự~\cite{groner2014vadl} mở rộng đúng kỹ thuật này sang
cả choreography, tức tương tác giữa nhiều tiến trình độc lập qua trao đổi thông điệp. Guizzardi và
cộng sự~\cite{guizzardi2014vemi} theo một hướng khác trong cùng nhóm: hợp nhất metamodel của i* và
metamodel của một nền tảng phát triển hướng mô hình thành một Integration Metamodel duy nhất, rồi
dùng các phép đo kiểm OCL thực thi được để phát hiện phần tử không thể ánh xạ và lớp không thể khởi
tạo, trước khi phép biến đổi mô hình thực sự diễn ra. Eshuis và Ghose~\cite{eshuis2021consistency} áp
dụng kỹ thuật kiểm tra nhất quán cấu trúc tương tự cho cặp mô hình mục tiêu và lược đồ quản lý case
khai báo. Li và cộng sự~\cite{li2015cim} đề xuất một chuỗi hình thức hoá khác hẳn về nền tảng toán
học, kết hợp lý thuyết phạm trù cho mô hình mục tiêu, mạng Petri cho mô hình tiến trình, và QVTo cho
các phép ánh xạ giữa hai mức.

Điểm mạnh chung của nhóm này là mức độ tự động hoá cao và cơ sở lý thuyết chặt: suy luận Description
Logics là decidable và bao phủ toàn bộ không gian trạng thái mà không cần liệt kê, trong khi việc
tái sử dụng công cụ OCL và EMF trưởng thành giúp Guizzardi và cộng sự~\cite{guizzardi2014vemi} phát
hiện lỗi trước khi biến đổi mô hình, khi chi phí sửa còn thấp. Tuy nhiên, phạm vi ngữ nghĩa của nhóm
này còn hẹp: kỹ thuật Description Logics chỉ xử lý quan hệ AND, IOR, và XOR decomposition, chưa mã
hoá contribution link, và gặp rủi ro undecidable khi quy trình có vòng lặp; công trình của Guizzardi
và cộng sự~\cite{guizzardi2014vemi} gắn chặt với một nền tảng phát triển hướng mô hình cụ thể và chỉ
kiểm tra được vi phạm cấu trúc, không phát hiện xung đột ở mức logic nghiệp vụ. Phương pháp đề xuất
khác biệt với nhóm này ở hai điểm. Thứ nhất, phương pháp đề xuất kiểm tra dọc theo một vết thực thi
cụ thể thay vì suy luận tĩnh tại một thời điểm, nên phát hiện được vi phạm phụ thuộc vào thứ tự thực
thi mà suy luận tĩnh bỏ sót. Thứ hai, phương pháp đề xuất đặc tả thêm được ràng buộc cấu trúc tổ
chức trong cùng miền OCL/USE dùng để kiểm i* và BPMN, điều không công trình nào trong nhóm này làm
được.

\subsection{Ngữ Nghĩa Vận Hành Theo Vết Thực Thi}\label{sec:rw-operational}

Nhóm công trình thứ hai định nghĩa một khái niệm vết hoặc trạng thái vận hành của tiến trình theo
thời gian, rồi đối chiếu vết đó với đặc tả logic hoặc thời gian của goal. Koliadis và
Ghose~\cite{koliadis2006goalbpm} đề xuất khung GoalBPM: gán mỗi hoạt động BPMN một effect annotation
ở thể tường thuật, trích xuất các trajectory từ sự kiện bắt đầu tới sự kiện kết thúc, rồi đối chiếu
hiệu ứng tích luỹ dọc mỗi trajectory với mẫu thời gian RT-LTL của goal KAOS, gồm Achieve, Maintain,
Avoid, và Cease. Dang và cộng sự~\cite{dang2010jucs} đề xuất khung khái niệm kịch bản, hành động, và
snapshot đã trình bày ở Phần~\ref{sec:introduction}, dùng triple graph grammar và OCL để đồng bộ hai
kịch bản ở hai mức trừu tượng khác nhau. Caballero-Villalobos và cộng
sự~\cite{caballero2026alig} dịch cả goal model iStar lẫn process model, workflow net hoặc DCR graph,
sang cùng một labeled transition system, dựng product của hai LTS đó thông qua một mapping function
nối process action với intentional element, rồi kiểm hai mức compliance lồng nhau bằng duyệt đồ thị
hai chiều trên product: weak compliance, khi mọi trạng thái đạt được đều có đường tiếp tục tới một
trạng thái thoả mọi quality, và strong compliance, tức weak compliance cộng thêm tính đơn điệu thoả
mãn quality theo suốt một run. Nagel và cộng sự~\cite{nagel2013consistency} kiểm tra sự phụ thuộc
logic và thời gian giữa các business goal khi chúng được ánh xạ sang một mô hình BPMN; một công trình
khác kiểm chứng bằng model checking hình thức trên một Requirement Dependency Graph mở rộng, đối
chiếu với mô hình BPMN cộng tác~\cite{n3_2019_reqdepgraph}.

So với nhóm suy luận tĩnh, nhóm này biểu diễn được thứ tự và tích luỹ nhân quả qua nhiều bước, và
một số công trình, như của Caballero-Villalobos và cộng sự~\cite{caballero2026alig}, có chứng minh
hình thức về tính đúng đắn và đầy đủ của thuật toán. Hạn chế của nhóm là khung GoalBPM của Koliadis
và Ghose~\cite{koliadis2006goalbpm} được thực hiện hoàn toàn thủ công; các kỹ thuật dựa trên hệ
chuyển trạng thái đối mặt bùng nổ không gian trạng thái khi quy trình có vòng lặp; và bản thân
Caballero-Villalobos và cộng sự~\cite{caballero2026alig} tự nhận rằng buộc mọi nhánh thực thi phải
dẫn tới trạng thái thoả mãn có thể tạo ra phán quyết vô lý về nghiệp vụ khi một số nhánh, như từ
chối một đơn xin phép, vốn dĩ không nên bị buộc phải đạt goal ban đầu. Phương pháp đề xuất khác biệt
với nhóm này ở hai điểm. Thứ nhất, snapshot pattern trong phương pháp đề xuất được suy tự động từ
OCL sẵn có trong mô hình, không cần viết tay như GoalBPM, và không cần tự xây luật lan truyền
AND/OR/Make/Break riêng như Caballero-Villalobos và cộng sự. Thứ hai, phương pháp đề xuất kiểm tra
đồng thời ràng buộc cấu trúc tổ chức trên cùng vết thực thi, không chỉ tính nhất quán goal--process.

\subsection{Căn Chỉnh Định Lượng Và Liên Tục}\label{sec:rw-quantitative}

Thay vì trả lời nhị phân đạt hay không đạt, nhóm công trình thứ ba tính một điểm số hoặc mức độ liên
tục. Guizzardi và Reis~\cite{guizzardi2012align} đề xuất phương pháp năm bước để căn chỉnh một mô
hình mục tiêu Tropos và một mô hình tiến trình BPMN đã tồn tại độc lập: phân loại execution path,
ánh xạ activity sang plan kèm trọng số đóng góp, phát hiện tự động goal chưa được hỗ trợ và activity
không đóng góp vào goal nào, rồi lan truyền điểm số thoả mãn theo từng path bằng quy tắc AND lấy giá
trị nhỏ nhất và OR lấy giá trị lớn nhất. Shamsaei và cộng sự~\cite{shamsaei2011kpi} theo dõi mức độ
tuân thủ quy trình bằng chỉ số hiệu suất KPI gắn với goal model GRL trong chuẩn URN. Một hướng ngược
lại tổng hợp mô hình mục tiêu từ các milestone của một mô hình tiến trình khai báo dữ
liệu~\cite{n9_2025_goalsynthesis}, cùng thuộc bản chất định lượng và liên tục dù đi theo chiều
ngược.

Điểm mạnh nổi bật của phương pháp năm bước~\cite{guizzardi2012align} là áp dụng được cho tình huống
cả hai mô hình đã tồn tại độc lập, không đòi hỏi quan hệ sinh một chiều như nhóm chuyển đổi. Hạn chế
chung của nhóm là việc gán trọng số mang tính chủ quan và phần lớn chưa được tự động hoá đầy đủ.
Phương pháp đề xuất khác biệt với nhóm này ở chỗ trả về ba danh sách lỗi tách theo chiều, cấu trúc,
quy trình, và goal, thay vì một điểm số tổng hợp duy nhất, và không yêu cầu chuyên gia gán trọng số
chủ quan cho từng liên kết đóng góp.

\subsection{Đối Chiếu Với Ràng Buộc Phẳng Bên Ngoài}\label{sec:rw-external}

Nhóm công trình thứ tư đối chiếu tiến trình với một đặc tả bên ngoài phẳng, không mang cấu trúc ý
định của goal model. Governatori và cộng sự~\cite{governatori2006compliance} hình thức hoá nghĩa vụ
hợp đồng và kiểm tra tiến trình có tuân thủ hợp đồng đó hay không. Ghose và
Koliadis~\cite{ghose2007auditing} kiểm toán tiến trình đối chiếu với quy định, một bài toán khác với
công trình của cùng hai tác giả đã khảo sát ở Phần~\ref{sec:rw-static}. Awad và
Weske~\cite{awad2009visualization} biểu diễn ràng buộc tuân thủ bằng mẫu thời gian LTL và trực quan
hoá vị trí vi phạm ngay trên mô hình BPMN. Van der Aalst và de Medeiros~\cite{aalst2005conformance}
đối chiếu trực tiếp log thực thi thật với một mô hình tiến trình để phát hiện thực thi bất thường.

Nhóm này thừa hưởng nhiều kỹ thuật kiểm chứng và trực quan hoá trưởng thành, đã ứng dụng công nghiệp
rộng rãi cho bài toán tuân thủ nói chung, nhưng phía yêu cầu luôn bị phẳng hoá thành luật, chỉ số,
hoặc log trước khi kiểm tra. Phương pháp đề xuất khác biệt với nhóm này ở chỗ giữ nguyên cấu trúc ý
định actor, goal, refinement, và contribution của mô hình i* làm phía yêu cầu, thay vì phẳng hoá nó
trước khi kiểm tra.

\subsection{Đặc Tả Và Kiểm Chứng Cấu Trúc Vai Trò Tổ Chức}\label{sec:rw-structure}

Nhóm công trình cuối cùng không thuộc phạm vi goal--process mà đặc tả và kiểm chứng cấu trúc tổ
chức. Hübner và cộng sự~\cite{hubner2002moise,hubner2007moiseplus} đề xuất mô hình tổ chức MOISE và
MOISE+ cho hệ đa tác tử, với đặc tả cấu trúc gồm role, group, và cardinality cho từng role trong
từng group. Ray và cộng sự~\cite{ray2004rbacuml} hình thức hoá các ràng buộc kiểm soát truy cập dựa
trên vai trò, gồm cả ràng buộc phân tách nhiệm vụ ngăn một chủ thể đồng thời giữ hai vai trò xung
đột, bằng OCL trên mô hình UML.

Điểm mạnh của nhóm này là có metamodel hình thức hoá tốt cho câu hỏi ai được giữ vai trò gì, đồng
thời với vai trò nào khác; công trình của Ray và cộng sự~\cite{ray2004rbacuml} còn cho thấy việc
dịch loại ràng buộc này sang OCL là khả thi và kiểm chứng được bằng công cụ UML/OCL trưởng thành.
Hạn chế mấu chốt là cả hai hướng nghiên cứu này dừng hoàn toàn ở mức đặc tả cấu trúc tổ chức tĩnh,
tách rời khỏi mô hình ý định và mô hình hành vi của cùng hệ thống. Phương pháp đề xuất khác biệt với
nhóm này ở chỗ không dừng lại ở đặc tả và kiểm chứng cấu trúc tổ chức một mình, mà tích hợp kết quả
đó vào cùng một hệ thống đối tượng dùng để kiểm tra tính nhất quán với mô hình mục tiêu và mô hình
tiến trình.

\subsection{Khoảng Trống Còn Lại}\label{sec:rw-gap}

Tổng hợp lại, chưa có công trình nào trong năm nhóm trên trả lời đồng thời RQ1 và RQ2 trên cùng một
mô hình hệ thống cụ thể: bốn nhóm đầu chỉ xét cặp goal--process theo bốn cơ chế khác nhau, còn nhóm
thứ năm chỉ xét cấu trúc tổ chức tách biệt. Phương pháp đề xuất ở Phần~\ref{sec:approach} đến
Phần~\ref{sec:checking} lấp khoảng trống này bằng cách kế thừa khung kịch bản--snapshot của Dang và
cộng sự~\cite{dang2010jucs} cho phần đồng bộ, và triết lý dịch sang một hình thức luận chung rồi
truy vấn của Guizzardi và cộng sự~\cite{guizzardi2014vemi} cho phần cấu trúc, đồng thời bổ sung
chiều tổ chức và kiểm tra dọc theo một vết thực thi cụ thể trên cùng một hệ thống đối tượng USE duy
nhất.
